% Part: incompleteness
% Chapter: incompleteness-provability
% Section: second-incompleteness-thm

\documentclass[../../../include/open-logic-section]{subfiles}

\begin{document}

% SEGMENT OLP-0319-S01

\olfileid{inc}{inp}{2in}

\olsection{இரண்டாம் முழுமையின்மைத் தேற்றம்}

$\Th{PA}$ தன் சொந்த முரணின்மையை நிறுவாது என்ற கூற்றை எவ்வாறு
வெளிப்படுத்தலாம்? $\Th{PA}$ முரணானது என்று சொல்வது,
$\Th{PA} \Proves \eq[0][1]$ என்று சொல்வதற்குச் சமம். ஆகவே,
முரணின்மைக் கூற்றாக $\OCon[\Th{PA}]$ என்பதை
!!{sentence}~$\lnot \OProv[\Th{PA}](\gn{\eq[0][1]})$ என்று
எடுத்துக்கொள்ளலாம்; பின்வரும் தேற்றம் நமக்குத் தேவையான முடிவைத் தருகிறது:

\begin{thm}
\ollabel{thm:second-incompleteness}
$\Th{PA}$ முரணற்றது என்று கருதுக. அப்போது $\Th{PA}$,
$\OCon[\Th{PA}]$ ஐ !!{derive} செய்யாது.
\end{thm}

இந்தத் தேற்றம் $\OCon[\Th{PA}]$ ஐ எவ்வாறு பிரதிநிதித்துவப்படுத்துகிறோம்
என்பதைப் பொறுத்தது என்பதை நினைவில் கொள்வது முக்கியம்; அதாவது,
$\OProv[\Th{PA}](y)$ இன் குறிப்பிட்ட பிரதிநிதித்துவத்தைப் பொறுத்தது.
$\OProv[\Th{PA}](y)$ ஐப் பிரதிநிதித்துவப்படுத்தும் அந்த வாய்பாடு மூன்று
!!{derivability} நிபந்தனைகளையும் நிறைவேற்றுகிறது என்பதையே நாம்
பயன்படுத்தப் போகிறோம். எனவே, இப்பண்புகளைக் கொண்ட
!!a{derivability} பயனிலை உள்ள எந்தக் கோட்பாட்டிற்கும் தேற்றம்
பொதுமைப்படுகிறது.

கோடலின் நிறுவல் வரைவைப் படிப்பது பயனுள்ளதாக இருக்கும்; ஏனெனில்
தேற்றம் ஒரு நல்ல நகைச்சுவையின் இறுதி வரியைப் போலத் திடீரென முடிகிறது.
அது பின்வருமாறு. \olref[1in]{thm:first-incompleteness} இன் நிறுவலில்
உருவாக்கிய கோடல் வாக்கியத்தை
$!G_\Th{PA}$ என்று கொள்வோம். “$\Th{PA}$ முரணற்றது எனில்,
$\Th{PA}$ ஆனது $!G_\Th{PA}$ ஐ !!{derive} செய்யாது” என்று நாம்
காட்டியுள்ளோம். இதை \emph{உள்ளேயே} $\Th{PA}$ இல் முறைப்படுத்தினால்,
\[
\OCon[\Th{PA}] \lif \lnot \OProv[\Th{PA}](\gn{!G_\Th{PA}}).
\]
என்ற நிறுவல் நமக்குக் கிடைக்கும். இப்போது $\Th{PA}$
$\OCon[\Th{PA}]$ ஐ !!{derive} செய்கிறது என்று கருதுக. அப்போது அது
$\lnot \Prov[\Th{PA}](\gn{!G_\Th{PA}})$ ஐயும் !!{derive} செய்யும்.
ஆனால் $!G_\Th{PA}$ ஒரு கோடல் வாக்கியம் என்பதால், இது
$!G_\Th{PA}$ க்கு நிகரானது. எனவே $\Th{PA}$,
$!G_\Th{PA}$ ஐ !!{derive} செய்யும்.

ஆனால், $\Th{PA}$ முரணற்றது எனில் அது $!G_\Th{PA}$ ஐ !!{derive}
செய்யாது என்பதை நாம் அறிவோம்! ஆகவே $\Th{PA}$ முரணற்றது எனில்,
அது $\OCon[\Th{PA}]$ ஐ !!{derive} செய்ய முடியாது.

வாதத்தை இன்னும் துல்லியமாக்க, $!G_\Th{PA}$ என்பதை $\Th{PA}$ க்கான
கோடல் வாக்கியமாக எடுத்துக்கொண்டு, (P1)--(P3) என்ற
!!{derivability} நிபந்தனைகளைப் பயன்படுத்தி
$\Th{PA}$, $\OCon[\Th{PA}] \lif !G_\Th{PA}$ ஐ !!{derive} செய்கிறது
என்று காட்டுவோம். இதனால் $\Th{PA}$,
$\OCon[\Th{PA}]$ ஐ !!{derive} செய்யாது என்பது தெரியும். இதோ
$\Th{PA}$ இல் நடைபெறும் நிறுவலின் வரைவு. (எளிமைக்காக,
$\Th{PA}$ என்ற கீழெழுத்துகளை விடுகிறோம்.)
\begin{align}
& !G \liff \lnot \OProv(\gn{!G}) \ollabel{G2-1}\\
& \qquad\text{$!G$ ஒரு கோடல் வாக்கியம்} \notag \\
& !G \lif \lnot \OProv(\gn{!G}) \ollabel{G2-2}\\
  & \qquad\text{\olref{G2-1} இலிருந்து} \notag\\
& !G \lif
  (\OProv(\gn{!G}) \lif \lfalse) \ollabel{G2-3}\\
  & \qquad\text{\olref{G2-2} இலிருந்து தருக்கத்தின் மூலம்} \notag\\
& \OProv(\gn{
    !G \lif
    (\OProv(\gn{!G}) \lif \lfalse)
  }) \ollabel{G2-4}\\
  & \qquad\text{\olref{G2-3} இலிருந்து P1 நிபந்தனையின் மூலம்} \notag\\
& \OProv(\gn{!G}) \lif
  \OProv(\gn{
    (\OProv(\gn{!G}) \lif \lfalse)
    }) \ollabel{G2-5}\\
  & \qquad\text{\olref{G2-4} இலிருந்து P2 நிபந்தனையின் மூலம்} \notag\\
& \OProv(\gn{!G}) \lif (\OProv(\gn{\OProv(\gn{!G})}) \lif \OProv(\gn{\lfalse})) \ollabel{G2-6}\\
  & \qquad\text{\olref{G2-5} இலிருந்து P2 நிபந்தனையும் தருக்கமும் மூலம்} \notag\\
& \OProv(\gn{!G}) \lif 
  \OProv(\gn{\OProv(\gn{!G})}) \ollabel{G2-7}\\
   & \qquad\text{P3 நிபந்தனையின் மூலம்} \notag\\
& \OProv(\gn{!G}) \lif \OProv(\gn{\lfalse}) \ollabel{G2-8}\\
  & \qquad \text{\olref{G2-6}, \olref{G2-7} ஆகியவற்றிலிருந்து தருக்கத்தின் மூலம்}\notag\\
& \OCon \lif \lnot \OProv(\gn{!G}) \ollabel{G2-9}\\
  & \qquad \text{\olref{G2-8} இன் எதிர்மாற்றமும்
  $\OCon \ident \lnot \OProv(\gn{\lfalse})$ என்பதும்} \notag \\
& \OCon \lif !G \notag\\
  & \qquad\text{\olref{G2-1}, \olref{G2-9} ஆகியவற்றிலிருந்து தருக்கத்தின் மூலம்} \notag
\end{align}
மேலுள்ள நிறுவலில் “தருக்கத்தின் மூலம்” என்று சொல்லப்படுவது
கூற்றுத் தருக்கத்தின் அடிப்படை உண்மைகளை மட்டுமே குறிக்கிறது.
எடுத்துக்காட்டாக, \olref{G2-3} இல்
$\Proves \lnot!A \liff (!A\lif\lfalse)$ என்பதையும்,
\olref{G2-8} இல்
$!A \lif (!B \lif !C), !A \lif !B \Proves !A \lif !C$
என்பதையும் பயன்படுத்துகிறோம். \olref{G2-5}, \olref{G2-6} ஆகியவற்றில்
P2 நிபந்தனையைப் பயன்படுத்துவது
$\OProv(\gn{!A \lif !B}) \lif
(\OProv(\gn{!A}) \lif \OProv(\gn{!B}))$ என்ற P2 இன் நிகழ்வுகளை
நம்பியுள்ளது. முதலாவது நிகழ்வில் $!A \ident !G$ என்றும்
$!B \ident \OProv(\gn{!G}) \lif \lfalse$ என்றும் எடுக்கிறோம்;
இரண்டாவதில் $!A \ident \OProv(\gn{G})$ என்றும்
$!B \ident \lfalse$ என்றும் எடுக்கிறோம்.

இரண்டாம் முழுமையின்மைத் தேற்றத்தின் இன்னும் பொதுவான வடிவம்
பின்வருமாறு:

\begin{thm}
\ollabel{thm:second-incompleteness-gen}
$\Th{Q}$ ஐ நீட்டிக்கும் எந்த முரணற்ற, !!{axiomatized} கோட்பாடான
$\Th{T}$ ஐயும் எடுத்துக்கொள்க. மேலும்
$\OProv[\Th{T}](y)$ என்பது $\Th{T}$ க்கான P1--P3
!!{derivability} நிபந்தனைகளை நிறைவேற்றும் ஏதேனும் வாய்பாடு என்க.
அப்போது $\Th{T}$, $\OCon[T]$ ஐ !!{derive} செய்யாது.
\end{thm}

\begin{prob}
$\Th{PA}$, $!G_{\Th{PA}} \lif \OCon[\Th{PA}]$ ஐ !!{derive} செய்கிறது
என்று காட்டுக.
\end{prob}

\begin{digress}
இதன் கருத்து என்னவென்றால், கணிதத்திற்கான எந்த “நியாயமான” முரணற்ற
கோட்பாடும் தன் சொந்த முரணின்மைக் கூற்றை !!{derive} செய்ய முடியாது.
$\Th{Q}$ ஐயும் ஹில்பெர்ட்டின் “இறுதிக்குட்பட்ட” காரணங்கூறலையும்
(அது எதுவாக இருந்தாலும்) உள்ளடக்கிய $\Th{T}$ என்ற கணிதக்
கோட்பாட்டைக் கருதுக. அப்போது முழு $\Th{T}$ உம் தன்
$\Th{T}$ முரணின்மைக்
கூற்றை !!{derive} செய்ய முடியாது; ஆகவே, அதைவிடவும், அதன்
இறுதிக்குட்பட்ட துண்டும் $\Th{T}$ இன் முரணின்மைக் கூற்றை !!{derive}
செய்ய முடியாது. அந்தப் பொருளில், “முழுக் கணிதத்திற்கும்”
இறுதிக்குட்பட்ட முரணின்மை நிறுவல் இருக்க முடியாது.

“இறுதிக்குட்பட்டது” என்ற சொல்லை விளக்குவதில் ஓரளவு நெகிழ்வு உண்டு.
1931 ஆம் ஆண்டுக் கட்டுரையில், நாம் “இறுதிக்குட்பட்டது” என்று
கருதக்கூடிய ஏதோ ஒன்று ஹில்பெர்ட் முறைப்படுத்த விரும்பிய கணிதத்தின்
வகைகளுக்கு வெளியே இருக்கலாம் என்ற வாய்ப்பை கோடல் ஏற்றுக்கொள்கிறார்.
ஆனால் கோடல் அன்புடன் இடமளித்தார்; இன்று, இறுதிக்குட்பட்டது என்று
நியாயமாக அழைக்கக்கூடியதும், தேர்வுக் கோட்பாட்டுடன் கூடிய
Zermelo--Fraenkel கணக் கோட்பாடான $\Th{ZFC}$ இல் முறைப்படுத்த
முடியாததுமான ஒன்றை எங்கே காண்போம் என்று புரிந்துகொள்வது கடினம்.
\end{digress}

\end{document}
